\documentclass[11pt,a4paper]{article}
\usepackage[spanish]{babel}
\usepackage[utf8]{inputenc}
\usepackage[T1]{fontenc}
\usepackage{lmodern}
\usepackage{inconsolata} % fuente monoespaciada profesional para código
\usepackage{geometry}
\geometry{margin=2cm}
\usepackage{microtype}
\usepackage{hyperref}
% Paquetes de navegación y estilo de índice
\usepackage{bookmark}
\usepackage{tocloft}
\usepackage{fancyhdr}
\usepackage{parskip}
\usepackage{listingsutf8}
\usepackage{upquote}
\usepackage[dvipsnames]{xcolor}
\usepackage{enumitem}
\usepackage{graphicx}
\usepackage{tikz}
\usepackage{tcolorbox}
\tcbuselibrary{skins,breakable, listings, theorems}
\usepackage{titlesec}
\usepackage{listings}
\usepackage{booktabs}
\usepackage{caption}
\captionsetup[lstlisting]{labelfont=bf, singlelinecheck=false, margin=0pt}

% Estilos visuales (paleta neutra profesional)
\definecolor{codebg}{HTML}{F5F7FA}
\definecolor{codefg}{HTML}{111827}
\definecolor{codekw}{HTML}{7C3AED}
\definecolor{codestr}{HTML}{8B4513}
\definecolor{codecm}{HTML}{6B7280}
\definecolor{codeid}{HTML}{111827}
\definecolor{accent}{HTML}{111827}
\definecolor{accent2}{HTML}{374151}

% Cajas estilizadas (con la nueva paleta)
\newtcolorbox{tipbox}{enhanced, breakable, colback=accent!5, colframe=accent!60!black, coltitle=black, title=Consejo Senior, attach boxed title to top left={yshift=-2mm, xshift=2mm}, boxed title style={colback=accent!15, sharp corners}}
\newtcolorbox{warnbox}{enhanced, breakable, colback=orange!5, colframe=orange!70!black, coltitle=black, title=Evita esto, attach boxed title to top left={yshift=-2mm, xshift=2mm}, boxed title style={colback=orange!20, sharp corners}}
\newtcolorbox{refbox}{enhanced, breakable, colback=green!5, colframe=green!60!black, coltitle=black, title=Patrón, attach boxed title to top left={yshift=-2mm, xshift=2mm}, boxed title style={colback=green!20, sharp corners}}

% Enlaces en negro
\hypersetup{colorlinks=true, linkcolor=black, urlcolor=black, citecolor=black}

% Índice (ToC) en español y con estilo
\addto\captionsspanish{\renewcommand{\contentsname}{Índice}}
\setcounter{tocdepth}{2}
\setcounter{secnumdepth}{3}
\renewcommand{\cfttoctitlefont}{\LARGE\bfseries}
\renewcommand{\cftaftertoctitle}{\par\vspace{6pt}\hrule\vspace{10pt}}
\renewcommand{\cftsecleader}{\cftdotfill{\cftdotsep}}
\renewcommand{\cftsecfont}{\bfseries}
\renewcommand{\cftsecpagefont}{\bfseries}
\setlength{\cftbeforesecskip}{4pt}
\setlength{\cftbeforesubsecskip}{2pt}

% Títulos con color y regla
\titleformat{\section}{\Large\bfseries\color{accent}}{\thesection}{0.6em}{}[\vspace{2pt}\titlerule]
\titlespacing*{\section}{0pt}{2.0ex plus 1ex minus .2ex}{1ex}
\titleformat{\subsection}{\large\bfseries\color{accent2}}{\thesubsection}{0.6em}{}
\titlespacing*{\subsection}{0pt}{1.6ex plus .8ex minus .2ex}{0.8ex}
\titleformat{\subsubsection}{\normalsize\bfseries}{\thesubsubsection}{0.6em}{}

% Encabezados/pies
\pagestyle{fancy}
\fancyhf{}
\fancyhead[L]{Spring Senior}
\fancyhead[R]{\nouppercase{\leftmark}}
\fancyfoot[C]{\thepage}
\renewcommand{\headrulewidth}{0.4pt}
\renewcommand{\footrulewidth}{0pt}
\setlength{\headheight}{15pt}

% Listings en UTF-8 con mapeo de acentos
\lstset{
  inputencoding=utf8,
  literate={á}{{\'a}}1 {é}{{\'e}}1 {í}{{\'\i}}1 {ó}{{\'o}}1 {ú}{{\'u}}1
           {Á}{{\'A}}1 {É}{{\'E}}1 {Í}{{\'I}}1 {Ó}{{\'O}}1 {Ú}{{\'U}}1
           {ñ}{{\~n}}1 {Ñ}{{\~N}}1
           {ü}{{\"u}}1 {Ü}{{\"U}}1
           {…}{{\ldots}}1
}

% Lenguajes para listings
\lstdefinelanguage{Java}{
  morekeywords={abstract,assert,boolean,break,byte,case,catch,char,class,const,continue,default,do,double,else,enum,extends,final,finally,float,for,if,goto,implements,import,instanceof,int,interface,long,native,new,package,private,protected,public,return,short,static,strictfp,super,switch,synchronized,this,throw,throws,transient,try,void,volatile,record,var},
  sensitive=true,
  morecomment=[l]{//},
  morecomment=[s]{/*}{*/},
  morestring=[b]", 
  morestring=[b]'
}

\lstdefinelanguage{YAML}{
  morekeywords={true,false,null,yes,no,on,off},
  sensitive=false,
  morecomment=[l]{\#}
}

\lstdefinelanguage{Bash}{
  morekeywords={echo,cd,ls,cat,grep,sed,awk,tail,head,export,docker,mvn,gradle,java,javac},
  sensitive=false,
  morecomment=[l]{\#}
}

% Estilos de listado (tema profesional)
\lstdefinestyle{basecodestyle}{
  backgroundcolor=\color{codebg},
  basicstyle=\ttfamily\small\color{codefg},
  keywordstyle=\color{codekw}\bfseries,
  stringstyle=\color{codestr},
  commentstyle=\color{codecm}\itshape,
  identifierstyle=\color{codeid},
  numberstyle=\scriptsize\color{codecm},
  numbers=left,
  stepnumber=1,
  numbersep=8pt,
  frame=single,
  rulecolor=\color{accent2},
  framesep=6pt,
  xleftmargin=12pt,
  xrightmargin=4pt,
  showstringspaces=false,
  breaklines=true,
  tabsize=2,
  upquote=true,
  captionpos=b
}

\lstdefinestyle{javastyle}{style=basecodestyle, language=Java}
\lstdefinestyle{yamlstyle}{style=basecodestyle, language=YAML}
\lstdefinestyle{bashstyle}{style=basecodestyle, language=Bash}

% Ajuste de título con guión largo (em-dash) para compatibilidad
\title{Spring Boot a Nivel Senior: Guía Práctica e Internals}
\author{Apuntes de trabajo}
\date{\today}

\begin{document}
\maketitle
\tableofcontents
\newpage

\section{Introducción}
Estos apuntes cubren Spring Boot 3.x (Spring Framework 6), Java 17+ y el stack moderno: auto-configuración, observabilidad (Micrometer + OpenTelemetry), AOT y native image, seguridad, datos, reactividad y despliegue. Incluye \textbf{internals} para comprender cómo funciona por debajo y tomar decisiones de arquitectura con criterio.

\begin{tipbox}
Mentalidad senior: simplicidad, límites claros, contratos estables (interfaces/DTOs), observabilidad desde el día uno, automatización y \emph{shift-left testing}.
\end{tipbox}

% --- NUEVO: Glosario y fundamentos para principiantes ---
\section{Glosario ultra básico (explicado simple)}
\begin{itemize}[leftmargin=*]
  \item \textbf{Arrancar (bootstrap)}: iniciar la app. En Java: método \texttt{main} que llama a \texttt{SpringApplication.run(...)}.
  \item \textbf{Spring Boot}: capas y utilidades que configuran Spring “solas” por convención (\emph{starters}, autoconfiguración).
  \item \textbf{Bean}: objeto gestionado por Spring (vida, dependencias). Se crea con \texttt{@Component}, \texttt{@Service}, \texttt{@Repository} o \texttt{@Bean}.
  \item \textbf{IoC/DI}: Inversión de Control/ Inyección de Dependencias. Spring crea e inyecta los beans donde se necesitan.
  \item \textbf{Controller/RestController}: clase que atiende peticiones HTTP y devuelve respuestas (JSON).
  \item \textbf{Propiedades}: configuración externa (\texttt{application.yml}). Puedes leerlas con \texttt{@Value} o \texttt{@ConfigurationProperties}.
  \item \textbf{JPA}: estándar para mapear objetos a tablas (ORM). Spring Data JPA te da repositorios listos.
  \item \textbf{Transacción}: bloque de trabajo todo-o-nada en BD, controlado con \texttt{@Transactional}.
  \item \textbf{Actuator}: endpoints de salud/metricas/info para operar la app.
  \item \textbf{WebFlux/Reactor}: modelo reactivo (no bloqueante) para alto número de conexiones.
  \item \textbf{AOP}: “envolver” métodos con lógica transversal (métricas, seguridad, transacciones).
\end{itemize}

\section{Fundamentos prácticos: de cero a tu primera API}
Objetivo: entender \emph{qué es cada cosa} y \emph{para qué sirve} con un ejemplo mínimo.

\subsection{Cómo arranca una app}
\begin{lstlisting}[style=javastyle, caption={Aplicación mínima}]
// src/main/java/com/example/App.java
package com.example;
import org.springframework.boot.autoconfigure.SpringBootApplication;
import org.springframework.boot.SpringApplication;

@SpringBootApplication // activa escaneo de componentes y autoconfiguración
public class App {
  public static void main(String[] args) { SpringApplication.run(App.class, args); }
}
\end{lstlisting}
\textbf{Qué pasa}: Spring busca clases anotadas (beans), crea el “contenedor” (IoC) y las conecta.

\subsection{Tu primer endpoint REST}
\begin{lstlisting}[style=javastyle, caption={Hola mundo HTTP}]
// src/main/java/com/example/web/HelloController.java
package com.example.web;
import org.springframework.web.bind.annotation.*;

@RestController
class HelloController {
  @GetMapping("/hello") String hello(@RequestParam(defaultValue="Mundo") String name) {
    return "Hola, " + name;
  }
}
\end{lstlisting}
\textbf{Qué pasa}: \texttt{@RestController} expone métodos HTTP. \texttt{@GetMapping} atiende \texttt{GET /hello}.

\subsection{Servicios e inyección de dependencias}
\begin{lstlisting}[style=javastyle, caption={Separar lógica en un servicio}] 
// src/main/java/com/example/core/GreeterService.java
package com.example.core;
import org.springframework.stereotype.Service;

@Service
public class GreeterService { public String greet(String name){ return "Hola, " + name; } }

// src/main/java/com/example/web/HelloController.java
@RestController
class HelloController {
  private final GreeterService greeter;
  HelloController(GreeterService greeter) { this.greeter = greeter; } // inyección por constructor
  @GetMapping("/hello") String hello(@RequestParam(defaultValue="Mundo") String name) {
    return greeter.greet(name);
  }
}
\end{lstlisting}
\textbf{Idea}: el controlador solo orquesta; la lógica vive en servicios.

\subsection{Propiedades y configuración}
\begin{lstlisting}[style=yamlstyle, caption={application.yml}]
spring.application.name: demo
app.title: "Demo App"
\end{lstlisting}
\begin{lstlisting}[style=javastyle, caption={ConfigurationProperties tipado}] 
// src/main/java/com/example/config/AppProps.java
@ConfigurationProperties(prefix="app")
public record AppProps(String title) {}

// Activar binding
a.@EnableConfigurationProperties(AppProps.class)
\end{lstlisting}
\textbf{Idea}: configura fuera del código y bindéalo a tipos seguros.

\subsection{Datos con JPA e H2 (memoria)}
\begin{lstlisting}[style=javastyle, caption={Entidad y repositorio}] 
@Entity class Note { @Id @GeneratedValue Long id; String text; }
interface NoteRepo extends JpaRepository<Note, Long> {}
\end{lstlisting}
\begin{lstlisting}[style=yamlstyle, caption={BD en memoria para desarrollo}]
spring.datasource.url: jdbc:h2:mem:demo;DB_CLOSE_DELAY=-1
spring.jpa.hibernate.ddl-auto: update
\end{lstlisting}
\textbf{Idea}: con pocas líneas guardas y lees datos reales.

\subsection{Validación y manejo de errores}
\begin{lstlisting}[style=javastyle, caption={Validar entrada y responder 400}] 
record CreateNote(@jakarta.validation.constraints.NotBlank String text) {}

@RestController class NotesCtrl {
  private final NoteRepo repo;
  NotesCtrl(NoteRepo repo){ this.repo = repo; }
  @PostMapping("/notes") Note create(@Valid @RequestBody CreateNote body){
    var n = new Note(); n.text = body.text(); return repo.save(n);
  }
}
\end{lstlisting}
\textbf{Idea}: anota DTOs y deja que Spring devuelva 400 automáticamente si no cumplen.

\begin{tipbox}
Construye siempre así: (1) endpoint mínimo, (2) servicio, (3) propiedades, (4) datos, (5) validación/errores.
\end{tipbox}

% --- FIN bloque de fundamentos ---

\section{Ecosistema y proyecto}
\paragraph{Qué significa}
Spring Boot trae “piezas” preconfiguradas (web, datos, seguridad, observabilidad). Arrancas rápido sin pelear con configuración de bajo nivel.
\paragraph{En la práctica}
Elige \emph{starters} según la necesidad (\texttt{spring-boot-starter-web}, \texttt{-data-jpa}, \texttt{-security}). Usa perfiles (\texttt{dev/prod}) en \texttt{application-*.yml}.

\begin{itemize}[leftmargin=*]
  \item \textbf{Spring Boot 3.x}: migración a \texttt{jakarta.*}, configuración por convención y \emph{starters}.\
  \item \textbf{Build}: Maven/Gradle, Buildpacks para imágenes, perfiles y entornos.\
  \item \textbf{Observabilidad}: Actuator, Micrometer, trazas con OTel.\
  \item \textbf{Ejecución}: JVM, contenedores, \emph{native images} con GraalVM.
\end{itemize}

\subsection{Estructura mínima}
\paragraph{Qué mirar}
La clase \texttt{@SpringBootApplication} y el \texttt{application.yml}. Luego, paquetes por \textbf{feature}: \texttt{web/}, \texttt{core/}, \texttt{data/}.

\begin{lstlisting}[style=javastyle, caption={Aplicación principal y configuración de beans}]
// src/main/java/com/example/App.java
package com.example;

import org.springframework.boot.SpringApplication;
import org.springframework.boot.autoconfigure.SpringBootApplication;
import org.springframework.context.annotation.Bean;
import java.util.concurrent.Executor;
import java.util.concurrent.Executors;

@SpringBootApplication
public class App {
  public static void main(String[] args) { SpringApplication.run(App.class, args); }

  // Executor con virtual threads (Java 21+, Spring Boot 3.2+)
  @Bean Executor appExecutor() { return Executors.newVirtualThreadPerTaskExecutor(); }
}
\end{lstlisting}

\begin{lstlisting}[style=yamlstyle, caption={application.yml con perfiles y observabilidad}]
spring:
  application:
    name: demo
  profiles:
    active: dev

management:
  endpoints:
    web:
      exposure:
        include: "*"
  tracing:
    enabled: true
    sampling:
      probability: 1.0
  metrics:
    tags:
      application: ${spring.application.name}
\end{lstlisting}

\section{Configuración tipada y propiedades}
\paragraph{Por qué importa}
Evita strings sueltos y errores de configuración. Con records y validación fallarás rápido y claro al arrancar.

\begin{lstlisting}[style=javastyle, caption={Configuration Properties con records y validación}]
// src/main/java/com/example/config/AppProps.java
package com.example.config;

import org.springframework.boot.context.properties.ConfigurationProperties;
import org.springframework.validation.annotation.Validated;
import jakarta.validation.constraints.NotBlank;
import java.time.Duration;

@Validated
@ConfigurationProperties(prefix = "app")
public record AppProps(@NotBlank String title, Duration timeout) {}
\end{lstlisting}

\begin{lstlisting}[style=javastyle, caption={Activación de binding tipado}]
// src/main/java/com/example/config/Config.java
package com.example.config;

import org.springframework.boot.context.properties.EnableConfigurationProperties;
import org.springframework.context.annotation.Configuration;

@Configuration
@EnableConfigurationProperties(AppProps.class)
class Config {}
\end{lstlisting}

\begin{lstlisting}[style=yamlstyle, caption={Valores de ejemplo}]
app:
  title: Demo App
  timeout: 5s
\end{lstlisting}

\section{Auto-configuración: cómo funciona por debajo}
\paragraph{Idea clave}
Spring Boot activa beans según “si está presente tal clase/propiedad y falta tal bean”. Así te ahorra wiring manual.
\paragraph{En la práctica}
Si algo no se activa, revisa el \emph{Condition Evaluation Report} o \texttt{/actuator/conditions}.

\begin{itemize}[leftmargin=*]
  \item Spring Boot importa \textbf{autoconfiguraciones} listadas en \texttt{META-INF/spring/org.springframework.boot.autoconfigure.AutoConfiguration.imports}.
  \item Cada clase \texttt{@AutoConfiguration} declara \texttt{@ConditionalOnClass}, \texttt{OnMissingBean}, \texttt{OnProperty}, etc. El \emph{ConditionEvaluationReport} explica por qué se activa o no.
  \item Orden: \texttt{@AutoConfigureBefore/After}. Desactivar con \texttt{spring.autoconfigure.exclude}.
\end{itemize}

\begin{lstlisting}[style=javastyle, caption={AutoConfiguration mínima condicional}]
package com.example.starter;

import org.springframework.boot.autoconfigure.AutoConfiguration;
import org.springframework.boot.autoconfigure.condition.ConditionalOnClass;
import org.springframework.boot.autoconfigure.condition.ConditionalOnMissingBean;
import org.springframework.context.annotation.Bean;

@AutoConfiguration
@ConditionalOnClass(name = "com.example.lib.Client")
public class ClientAutoConfiguration {

  @Bean
  @ConditionalOnMissingBean
  public ClientApi clientApi() { return new ClientApi(); }
}
\end{lstlisting}

\begin{lstlisting}[style=yamlstyle, caption={Registrar la autoconfiguración}]
# META-INF/spring/org.springframework.boot.autoconfigure.AutoConfiguration.imports
com.example.starter.ClientAutoConfiguration
\end{lstlisting}

\begin{tipbox}
Para depurar condiciones: arranca con \texttt{--debug} o revisa el \emph{Condition Evaluation Report} en logs. Usa Actuator endpoint \texttt{/actuator/conditions}.
\end{tipbox}

\section{Web MVC y REST}
\paragraph{Cómo diseñar}
Controladores finos que validan y orquestan; servicios con la lógica; DTOs para entrada/salida; \texttt{@ControllerAdvice} para errores uniformes.

\begin{lstlisting}[style=javastyle, caption={Controlador REST con validación y manejo de errores}]
package com.example.web;

import org.springframework.web.bind.annotation.*;
import jakarta.validation.Valid;
import jakarta.validation.constraints.*;
import org.springframework.http.ResponseEntity;

record UserDto(@NotBlank String id, @Email String email) {}

@RestController
@RequestMapping("/api/users")
class UserController {
  @GetMapping("/{id}")
  ResponseEntity<UserDto> get(@PathVariable String id) {
    return ResponseEntity.ok(new UserDto(id, "user@example.com"));
  }

  @PostMapping
  ResponseEntity<UserDto> create(@Valid @RequestBody UserDto body) {
    return ResponseEntity.status(201).body(body);
  }
}
\end{lstlisting}

\begin{lstlisting}[style=javastyle, caption={Advice global para errores}]
package com.example.web;

import org.springframework.web.bind.annotation.*;
import org.springframework.http.*;
import org.springframework.web.method.annotation.MethodArgumentTypeMismatchException;
import org.springframework.validation.BindException;

@RestControllerAdvice
class GlobalErrors {
  @ExceptionHandler({ BindException.class, MethodArgumentTypeMismatchException.class })
  ResponseEntity<?> badReq(Exception ex) {
    return ResponseEntity.badRequest().body(new ErrorPayload("BAD_REQUEST", ex.getMessage()));
  }

  record ErrorPayload(String code, String message) {}
}
\end{lstlisting}

\section{Datos: JPA, transacciones y migraciones}
\paragraph{Reglas prácticas}
- Transacciones en servicios. - Evita exponer entidades en la API, usa DTOs. - Controla N+1 con \texttt{EntityGraph} o \emph{fetch join}.

\begin{lstlisting}[style=javastyle, caption={Entidad JPA y repositorio}] 
package com.example.data;

import jakarta.persistence.*;
import org.springframework.data.jpa.repository.*;
import org.springframework.data.repository.query.Param;
import java.util.*;

@Entity
@Table(name = "users")
public class User {
  @Id @GeneratedValue(strategy = GenerationType.IDENTITY)
  Long id;
  @Column(nullable = false, unique = true)
  String email;
}

interface UserRepo extends JpaRepository<User, Long> {
  Optional<User> findByEmail(@Param("email") String email);
}
\end{lstlisting}

\begin{lstlisting}[style=javastyle, caption={Servicio con @Transactional y reglas clave}]
package com.example.data;

import org.springframework.stereotype.Service;
import org.springframework.transaction.annotation.Transactional;

@Service
class UserService {
  private final UserRepo repo;
  UserService(UserRepo repo) { this.repo = repo; }

  @Transactional
  public User create(String email) {
    var u = new User();
    u.email = email;
    return repo.save(u);
  }
}
\end{lstlisting}

\begin{lstlisting}[style=yamlstyle, caption={Configuración JPA y Flyway}]
spring:
  datasource:
    url: jdbc:postgresql://localhost:5432/app
    username: app
    password: secret
  jpa:
    hibernate:
      ddl-auto: validate
    properties:
      hibernate:
        format_sql: true
        jdbc.batch_size: 50
  flyway:
    enabled: true
    locations: classpath:db/migration
\end{lstlisting}

\begin{refbox}
Evita LazyInitializationException: diseña límites de transacción en servicios, usa \emph{fetch joins} o \emph{entity graphs}. Controla el N+1 con logs y \texttt{hibernate.default\_batch\_fetch\_size}.
\end{refbox}

\section{Seguridad con Spring Security 6}
\paragraph{Modelo mental}
Una cadena de filtros decide quién eres (autenticación) y qué puedes hacer (autorización). Mantén reglas simples y centralizadas.

\begin{lstlisting}[style=javastyle, caption={SecurityFilterChain, CORS y JWT Resource Server}]
package com.example.security;

import org.springframework.context.annotation.*;
import org.springframework.security.config.annotation.web.builders.HttpSecurity;
import org.springframework.security.web.SecurityFilterChain;
import org.springframework.security.config.Customizer;

@Configuration
class SecurityConfig {
  @Bean
  SecurityFilterChain http(HttpSecurity http) throws Exception {
    http
      .csrf(csrf -> csrf.disable())
      .cors(Customizer.withDefaults())
      .authorizeHttpRequests(auth -> auth
        .requestMatchers("/actuator/**").permitAll()
        .requestMatchers("/api/admin/**").hasRole("ADMIN")
        .anyRequest().authenticated())
      .oauth2ResourceServer(oauth -> oauth.jwt(Customizer.withDefaults()));
    return http.build();
  }
}
\end{lstlisting}

\begin{lstlisting}[style=yamlstyle, caption={Propiedades JWT (ejemplo con jwk-set-uri)}]
spring:
  security:
    oauth2:
      resourceserver:
        jwt:
          jwk-set-uri: https://issuer.example.com/.well-known/jwks.json
\end{lstlisting}

\begin{tipbox}
Método seguro: validación a nivel de método con \texttt{@EnableMethodSecurity} y anotaciones \texttt{@PreAuthorize("hasRole('ADMIN')")}. Externaliza autorizaciones complejas a un servicio.
\end{tipbox}

\section{Reactividad y WebFlux}
\paragraph{Cuándo usarlo}
I/O intensivo con muchas conexiones y latencia de redes. Si tu stack es bloqueante (JPA), quédate en MVC.

\begin{lstlisting}[style=javastyle, caption={Controlador reactivo y WebClient}]
package com.example.reactive;

import org.springframework.web.bind.annotation.*;
import reactor.core.publisher.*;
import org.springframework.web.reactive.function.client.WebClient;

@RestController
class ReactiveCtrl {
  private final WebClient http = WebClient.create();

  @GetMapping("/rx") Mono<String> hello() {
    return Mono.just("ok");
  }

  @GetMapping("/rx/users") Flux<String> users() {
    return http.get().uri("https://example.com/users")
      .retrieve().bodyToFlux(String.class);
  }
}
\end{lstlisting}

\begin{warnbox}
No mezcles WebFlux con JPA bloqueante; usa R2DBC o drivers no bloqueantes. Si debes integrar, ejecuta en \texttt{boundedElastic} pero entiende el coste.
\end{warnbox}

\section{Mensajería, tareas y caché}
\subsection{Kafka / RabbitMQ}
\begin{lstlisting}[style=javastyle, caption={Productor/consumidor Kafka}]
package com.example.msg;

import org.springframework.kafka.annotation.KafkaListener;
import org.springframework.kafka.core.KafkaTemplate;
import org.springframework.stereotype.Service;

@Service
class KafkaSvc {
  private final KafkaTemplate<String, String> kt;
  KafkaSvc(KafkaTemplate<String, String> kt) { this.kt = kt; }
  void publish(String topic, String value) { kt.send(topic, value); }

  @KafkaListener(topics = "events", groupId = "app")
  void consume(String value) { System.out.println("event=" + value); }
}
\end{lstlisting}

\subsection{Tareas y scheduling}
\begin{lstlisting}[style=javastyle, caption={@Async y @Scheduled con pools dedicados}]
package com.example.bg;

import org.springframework.scheduling.annotation.*;
import org.springframework.context.annotation.*;
import java.util.concurrent.*;

@Configuration
@EnableAsync
@EnableScheduling
class BgConfig {
  @Bean Executor taskExecutor() { return Executors.newFixedThreadPool(8); }

  @Async
  public void work() { /* ... */ }

  @Scheduled(cron = "0 */5 * * * *")
  public void job() { /* ... */ }
}
\end{lstlisting}

\subsection{Caché}
\begin{lstlisting}[style=javastyle, caption={Caché con Caffeine/Redis}]
package com.example.cache;

import org.springframework.cache.annotation.*;
import org.springframework.context.annotation.Configuration;

@Configuration
@EnableCaching
class CacheCfg {}

class ProductService {
  @Cacheable(cacheNames = "products", key = "#id")
  Product findById(Long id) { /* lookup */ return null; }
}
\end{lstlisting}

\section{Observabilidad con Actuator y Micrometer}
\paragraph{Qué exponer}
Salud, métricas y trazas. Etiqueta métricas con nombre de aplicación y versión. Envía trazas a tu collector (OTel).

\begin{lstlisting}[style=yamlstyle, caption={Actuator y exportación de métricas/trazas}]
management:
  endpoints.web.exposure.include: "*"
  metrics.export.prometheus.enabled: true
  otlp:
    tracing:
      endpoint: http://otel-collector:4318/v1/traces
\end{lstlisting}

\begin{lstlisting}[style=javastyle, caption={Métricas personalizadas}]
package com.example.obs;

import io.micrometer.core.instrument.MeterRegistry;
import org.springframework.stereotype.Component;

@Component
class Metrics {
  Metrics(MeterRegistry registry) {
    registry.counter("app.startups").increment();
  }
}
\end{lstlisting}

\section{Testing eficaz}
\paragraph{Estrategia mínima}
- Unit: servicios con mocks. - Slices: \texttt{@WebMvcTest}, \texttt{@DataJpaTest}. - Integración: Testcontainers para BD real. - Smoke: endpoints de Actuator.

\begin{lstlisting}[style=javastyle, caption={@SpringBootTest + MockMvc}]
package com.example;

import org.junit.jupiter.api.Test;
import org.springframework.boot.test.context.SpringBootTest;
import org.springframework.boot.test.autoconfigure.web.servlet.AutoConfigureMockMvc;
import org.springframework.test.web.servlet.MockMvc;
import org.springframework.beans.factory.annotation.Autowired;

import static org.springframework.test.web.servlet.request.MockMvcRequestBuilders.get;
import static org.springframework.test.web.servlet.result.MockMvcResultMatchers.status;

@SpringBootTest
@AutoConfigureMockMvc
class AppTests {
  @Autowired MockMvc mvc;

  @Test void health() throws Exception {
    mvc.perform(get("/actuator/health")).andExpect(status().isOk());
  }
}
\end{lstlisting}

\begin{lstlisting}[style=javastyle, caption={Testcontainers para PostgreSQL}]
package com.example;

import org.junit.jupiter.api.*;
import org.springframework.boot.test.context.SpringBootTest;
import org.testcontainers.containers.PostgreSQLContainer;
import org.springframework.test.context.DynamicPropertyRegistry;
import org.springframework.test.context.DynamicPropertySource;

@SpringBootTest
class PgIT {
  static PostgreSQLContainer<?> pg = new PostgreSQLContainer<>("postgres:15");
  static { pg.start(); }

  @DynamicPropertySource
  static void props(DynamicPropertyRegistry r) {
    r.add("spring.datasource.url", pg::getJdbcUrl);
    r.add("spring.datasource.username", pg::getUsername);
    r.add("spring.datasource.password", pg::getPassword);
  }
}
\end{lstlisting}

\section{Empaquetado, Docker y Native Image}
\paragraph{Ruta práctica}
Desarrollo con \texttt{spring-boot:run}, empaqueta JAR, luego imagen con Buildpacks. Prueba Native si necesitas arranque ultrarrápido y menor memoria.

\begin{lstlisting}[style=bashstyle, caption={Comandos útiles Maven/Gradle y contenedores}]
# Ejecutar en desarrollo
mvn spring-boot:run

# Empaquetar jar
mvn -DskipTests package

# Imagen con buildpacks
mvn spring-boot:build-image -DskipTests

# Imagen nativa (GraalVM Native Build Tools instalado)
mvn -Pnative native:compile
\end{lstlisting}

\begin{warnbox}
AOT/native reduce memoria y latencia de arranque pero limita reflexión y proxys dinámicos. Verifica compatibilidad de librerías (hints) y usa \texttt{@RegisterReflectionForBinding} cuando sea necesario.
\end{warnbox}

\section{Rendimiento y escalabilidad}
\paragraph{Por dónde empezar}
Mide primero. Ajusta pool de conexiones, usa paginación, evita trabajo en el hilo de petición, añade caché donde tenga sentido.

\begin{itemize}[leftmargin=*]
  \item Conexiones: HikariCP, tamaño de pool acorde al workload; en reactividad, evita pools bloqueantes.
  \item GC y memoria: \texttt{-Xms/-Xmx}, G1/Generational ZGC según JDK.
  \item Virtual Threads: alto \emph{throughput} I/O-bound con código imperativo; mide latencias y bloqueo de drivers.
  \item Warmup y JIT: \emph{tiered compilation}; en nativo no hay JIT, calibra tiempo de arranque vs rendimiento sostenido.
\end{itemize}

\section{Patrones y anti-patrones}
\begin{refbox}
Patrones: Hexagonal/Ports-Adapters, Domain Services, DTO/assembler, Outbox para consistencia entre BDs y colas, Idempotency keys.
\end{refbox}

\begin{warnbox}
Anti-patrones: entidades anémicas con lógica en controladores, transacciones cruzando límites externos, \texttt{@Transactional} en métodos \texttt{private}, \emph{God Services}, repositorios chatos que exponen JPA al resto del dominio.
\end{warnbox}

\section{Checklist de revisión senior}
\begin{enumerate}[leftmargin=*]
  \item Configuración tipada (\texttt{@ConfigurationProperties}) y validada.
  \item Límites de transacción claros; sin N+1; migraciones con Flyway/Liquibase.
  \item Seguridad centralizada; endpoints expuestos y CORS definidos.
  \item Observabilidad: Actuator, métricas y trazas; logs estructurados.
  \item Tests con slices y Testcontainers para integrar con dependencias.
  \item Dockerfile o Buildpacks; perfiles por entorno; \emph{readiness/liveness} para K8s.
\end{enumerate}

% ------------------------------------------------------
% SPRING CORE (no solo Boot): IoC, ciclo de vida, AOP, transacciones
% ------------------------------------------------------
\section{Spring Core avanzado (más allá de Spring Boot)}
\subsection{Contenedor IoC, ciclos de vida y extensibilidad}
\begin{itemize}[leftmargin=*]
  \item \textbf{BeanFactory vs ApplicationContext}: este último añade internacionalización, eventos, recursos y autowiring avanzado.
  \item \textbf{Ciclo de vida}: instanciación \textrightarrow{} inyección \textrightarrow{} callbacks (\texttt{InitializingBean}/\texttt{@PostConstruct}) \textrightarrow{} post-procesadores \textrightarrow{} \texttt{SmartInitializingSingleton} \textrightarrow{} destrucción (\texttt{DisposableBean}/\texttt{@PreDestroy}).
  \item \textbf{Scopes}: \texttt{singleton}, \texttt{prototype} y web (\texttt{request}/\texttt{session}). Cuidado con inyectar \texttt{prototype} en \texttt{singleton}.
  \item Extensión: \texttt{BeanFactoryPostProcessor}, \texttt{BeanDefinitionRegistryPostProcessor}, \texttt{BeanPostProcessor}, \texttt{FactoryBean}, eventos (\texttt{ApplicationEventPublisher}/\texttt{@EventListener}).
\end{itemize}

\begin{lstlisting}[style=javastyle, caption={BeanPostProcessor y FactoryBean}] 
// Post-procesador que aplica un wrapper condicional
@Component
class TimingPostProcessor implements BeanPostProcessor {
  @Override public Object postProcessAfterInitialization(Object bean, String name) {
    if (bean instanceof Service) return ProxyFactory.getProxy(bean.getClass(), new TimingInterceptor(bean));
    return bean;
  }
}

// FactoryBean crea un bean complejo bajo un id simple
@Component("client")
class ClientFactory implements FactoryBean<Client> {
  public Client getObject() { return new Client("cfg"); }
  public Class<?> getObjectType() { return Client.class; }
}
\end{lstlisting}

\begin{lstlisting}[style=javastyle, caption={Eventos y listeners transaccionales}] 
@Component
class Eventing {
  private final ApplicationEventPublisher pub;
  Eventing(ApplicationEventPublisher pub) { this.pub = pub; }
  public void fire(User u) { pub.publishEvent(new UserCreated(u.id())); }
}

record UserCreated(Long id) {}

@Component
class Audits {
  @EventListener
  void on(UserCreated e) { /* antes de commit */ }

  @TransactionalEventListener(phase = TransactionPhase.AFTER_COMMIT)
  void onCommitted(UserCreated e) { /* side-effects confiables */ }
}
\end{lstlisting}

\subsection{AOP en profundidad}
\begin{itemize}[leftmargin=*]
  \item \textbf{Proxies} JDK (interfaces) y CGLIB (clases). Problema clásico: \emph{self-invocation} no pasa por el proxy.
  \item Tipos de advice: \texttt{@Before}, \texttt{@AfterReturning}, \texttt{@AfterThrowing}, \texttt{@Around}. Diseño de pointcuts con AspectJ (\texttt{execution}, \texttt{within}, \texttt{@annotation}, etc.).
  \item Ordenación con \texttt{@Order}; minimizar \emph{pointcuts} amplios para evitar impacto en rendimiento.
\end{itemize}

\begin{lstlisting}[style=javastyle, caption={Aspecto de métricas con @Around}]
@Aspect
@Component
@Order(10)
class MetricsAspect {
  private final MeterRegistry m;
  MetricsAspect(MeterRegistry m) { this.m = m; }

  @Around("within(@org.springframework.stereotype.Service *)")
  Object timed(ProceedingJoinPoint pjp) throws Throwable {
    var sample = Timer.start(m);
    try { return pjp.proceed(); }
    finally { sample.stop(m.timer("svc.time", "method", pjp.getSignature().getName())); }
  }
}
\end{lstlisting}

\subsection{Transacciones avanzadas}
\begin{itemize}[leftmargin=*]
  \item Propagación: \texttt{REQUIRED} (por defecto), \texttt{REQUIRES\_NEW}, \texttt{NESTED}, \texttt{MANDATORY}, \texttt{SUPPORTS}, \texttt{NOT\_SUPPORTED}, \texttt{NEVER}.
  \item Aislamiento: \texttt{READ\_COMMITTED}, \texttt{REPEATABLE\_READ}, \texttt{SERIALIZABLE}, etc. \texttt{readOnly} es una \emph{pista}; con JPA puede optimizar el \emph{flush}.
  \item Reglas de rollback: por defecto para \emph{unchecked}; usa \texttt{rollbackFor} para \emph{checked}.
  \item Eventos transaccionales: \texttt{TransactionSynchronizationManager}; \texttt{@TransactionalEventListener} para \emph{after-commit}.
\end{itemize}

\begin{lstlisting}[style=javastyle, caption={Propagación y locking}]
@Service
class OrderService {
  private final Repo r;
  OrderService(Repo r) { this.r = r; }

  @Transactional
  public void place(Order o) { /* REQUIRED */ save(o); charge(o); }

  @Transactional(propagation = Propagation.REQUIRES_NEW)
  public void audit(Order o) { /* registra en otra tx */ }

  @Transactional(isolation = Isolation.REPEATABLE_READ)
  public Order findWithLock(Long id) {
    return r.findByIdForUpdate(id); // SELECT ... FOR UPDATE (pessimistic)
  }
}
\end{lstlisting}

\begin{warnbox}
No llames métodos \texttt{@Transactional} de la misma instancia (\emph{self-invocation}): el proxy no intercepta. Extrae a otro bean.
\end{warnbox}

% ------------------------------------------------------
% Spring Data avanzado
% ------------------------------------------------------
\section{Spring Data avanzado}
\begin{itemize}[leftmargin=*]
  \item Ciclo de vida JPA: estados (transient, managed, detached, removed), \emph{dirty checking} y \texttt{flush}.
  \item \textbf{Proyecciones} (interfaces/clases), \textbf{Specifications} y Querydsl; \textbf{EntityGraph} para asociaciones.
  \item Auditoría (\texttt{@CreatedDate}/\texttt{@LastModifiedDate}) y \texttt{AuditorAware}.
  \item Repositorios personalizados y \texttt{@Query} para DTOs.
\end{itemize}

\begin{lstlisting}[style=javastyle, caption={EntityGraph y proyecciones}] 
interface UserView { Long getId(); String getEmail(); }

interface UserRepo extends JpaRepository<User, Long>, JpaSpecificationExecutor<User> {
  @EntityGraph(attributePaths = {"roles"})
  Optional<User> findWithRolesById(Long id);

  @Query("select u.id as id, u.email as email from User u where u.email like :q")
  List<UserView> search(@Param("q") String q);
}
\end{lstlisting}

\begin{lstlisting}[style=javastyle, caption={Specifications para filtros dinámicos}] 
class Specs {
  static Specification<User> emailContains(String q) {
    return (root, cq, cb) -> cb.like(root.get("email"), "%" + q + "%");
  }
  static Specification<User> hasRole(String role) {
    return (root, cq, cb) -> cb.isMember(role, root.get("roles"));
  }
}
\end{lstlisting}

% ------------------------------------------------------
% Spring Security avanzado
% ------------------------------------------------------
\section{Spring Security avanzado}
\begin{itemize}[leftmargin=*]
  \item Arquitectura: cadena de filtros, \texttt{AuthenticationManager}, \texttt{AuthenticationProvider}, \texttt{SecurityContext}.
  \item Autorización moderna con \texttt{AuthorizationManager}; seguridad a nivel de método con \texttt{@EnableMethodSecurity}.
  \item Múltiples cadenas \texttt{HttpSecurity} con \texttt{@Order} para separar APIs públicas/privadas.
  \item \textbf{PermissionEvaluator} para reglas de dominio; sesiones \emph{stateless} y \texttt{SessionCreationPolicy.STATELESS}.
\end{itemize}

\begin{lstlisting}[style=javastyle, caption={Múltiples HttpSecurity y PermissionEvaluator}] 
@Configuration
@EnableMethodSecurity
class MultiHttp {
  @Bean @Order(1)
  SecurityFilterChain api(HttpSecurity http) throws Exception {
    http.securityMatcher("/api/**")
      .authorizeHttpRequests(a -> a.anyRequest().authenticated())
      .oauth2ResourceServer(o -> o.jwt());
    return http.build();
  }
  @Bean @Order(2)
  SecurityFilterChain web(HttpSecurity http) throws Exception {
    http.authorizeHttpRequests(a -> a.anyRequest().permitAll());
    return http.build();
  }
}

@Component
class DomainPermission implements PermissionEvaluator {
  public boolean hasPermission(Authentication a, Object target, Object perm) {
    return /* lógica de dominio */ true; }
  public boolean hasPermission(Authentication a, Serializable id, String type, Object perm) { return true; }
}
\end{lstlisting}

% ------------------------------------------------------
% Web MVC y WebFlux internals
% ------------------------------------------------------
\section{MVC/WebFlux: resolvers, handlers y funcional}
\subsection{Spring MVC}
\begin{itemize}[leftmargin=*]
  \item \textbf{ArgumentResolvers}: inyectan parámetros en controladores; \textbf{ReturnValueHandlers}: serializan respuestas.
  \item Negotiación de contenido, validación y \texttt{HttpMessageConverters} (Jackson, etc.).
\end{itemize}

\begin{lstlisting}[style=javastyle, caption={ArgumentResolver personalizado}]
@Configuration
class MvcCfg implements WebMvcConfigurer {
  @Override public void addArgumentResolvers(List<HandlerMethodArgumentResolver> resolvers) {
    resolvers.add(new UserIdResolver());
  }
}

class UserId {}
class UserIdResolver implements HandlerMethodArgumentResolver {
  public boolean supportsParameter(MethodParameter p) { return p.getParameterType().equals(UserId.class); }
  public Object resolveArgument(MethodParameter p, ModelAndViewContainer m, NativeWebRequest w, WebDataBinderFactory b) {
    return new UserId();
  }
}
\end{lstlisting}

\subsection{Spring WebFlux}
\begin{itemize}[leftmargin=*]
  \item Modelo \emph{reactivo} basado en \textbf{Reactor}: \emph{event loop}, \emph{backpressure} y \texttt{Schedulers}.
  \item Endpoints funcionales con \texttt{RouterFunction}/\texttt{HandlerFunction}.
\end{itemize}

\begin{lstlisting}[style=javastyle, caption={Router funcional en WebFlux}] 
@Configuration
class Routes {
  @Bean RouterFunction<ServerResponse> router() {
    return RouterFunctions.route()
      .GET("/hello", req -> ServerResponse.ok().bodyValue("hi"))
      .build();
  }
}
\end{lstlisting}

% ------------------------------------------------------
% Integration, Batch, Cloud, Resilience
% ------------------------------------------------------
\section{Integración, Batch y resiliencia}
\subsection{Spring Integration}
\begin{itemize}[leftmargin=*]
  \item Colas, canales, \emph{gateways}, transformadores y \emph{Enterprise Integration Patterns}.
\end{itemize}

\subsection{Spring Batch}
\begin{itemize}[leftmargin=*]
  \item \textbf{Jobs/Steps}: \emph{chunk-oriented} vs \emph{tasklet}. Reintentables, idempotencia y control de estado.
\end{itemize}

\subsection{Resilience4j y Spring Cloud}
\begin{itemize}[leftmargin=*]
  \item Circuit Breaker, Retry, RateLimiter, Bulkhead. Configuración por instancia.
  \item Spring Cloud: Config, Gateway, OpenFeign, LoadBalancer; trazas con Micrometer Tracing/OTel.
\end{itemize}

\begin{lstlisting}[style=yamlstyle, caption={Resilience4j por instancia}] 
resilience4j.circuitbreaker:
  instances:
    externalApi:
      slidingWindowSize: 20
      failureRateThreshold: 50
\end{lstlisting}

% ------------------------------------------------------
% Configuración avanzada
% ------------------------------------------------------
\section{Configuración avanzada y PropertySources}
\begin{itemize}[leftmargin=*]
  \item Orden de resolución: args de línea de comandos \textrightarrow{} variables de entorno \textrightarrow{} archivos \texttt{application-*.yml} \textrightarrow{} \texttt{@PropertySource}.
  \item \texttt{Environment} y \texttt{PropertySources}; perfiles y \texttt{@Profile}, \texttt{@Conditional} personalizadas.
  \item Cifrado de propiedades (JCE/HashiCorp Vault/SOPS) y secretos externos.
\end{itemize}

\begin{lstlisting}[style=javastyle, caption={Condition personalizada}]
class OnLinux implements Condition {
  public boolean matches(ConditionContext c, AnnotatedTypeMetadata m) {
    return c.getEnvironment().getProperty("os.name", "").toLowerCase().contains("linux");
  }
}

@Configuration
@Conditional(OnLinux.class)
class LinuxOnly {}
\end{lstlisting}

% ------------------------------------------------------
% Testing avanzado
% ------------------------------------------------------
\section{Testing avanzado}
\begin{itemize}[leftmargin=*]
  \item \textbf{Slices}: \texttt{@WebMvcTest}, \texttt{@DataJpaTest}, \texttt{@JsonTest}, \texttt{@RestClientTest}.
  \item \texttt{@MockBean} para sustituir dependencias; \textbf{context caching} para acelerar.
  \item WireMock/Testcontainers para dependencias externas (HTTP/Kafka/Redis).
\end{itemize}

\begin{lstlisting}[style=javastyle, caption={@WebMvcTest con servicio simulado}]
@WebMvcTest(UserController.class)
class UserControllerTest {
  @Autowired MockMvc mvc;
  @MockBean UserService svc;

  @Test void ok() throws Exception {
    when(svc.get(anyString())).thenReturn(new UserDto("1","a@b.com"));
    mvc.perform(get("/api/users/1")).andExpect(status().isOk());
  }
}
\end{lstlisting}

% ------------------------------------------------------
% Plan de estudio senior Spring
% ------------------------------------------------------
\section{Plan de estudio Senior Spring (6 semanas)}
\begin{enumerate}[leftmargin=*]
  \item Core IoC, ciclo de vida, eventos, AOP básico; labs con \texttt{BeanPostProcessor} y \texttt{FactoryBean}.
  \item Transacciones avanzadas, locking, specs JPA y proyecciones; pruebas con \texttt{@DataJpaTest} y Testcontainers.
  \item Security avanzado (multi-HttpSecurity, método), MVC internals (resolvers/handlers).
  \item WebFlux: router funcional, backpressure, integración con R2DBC.
  \item Integración/Batch/Resilience4j; observabilidad profunda con métricas y trazas.
  \item Cloud: Config, Gateway, OpenFeign; despliegue en contenedores y rendimiento.
\end{enumerate}

% ------------------------------------------------------
% API y serialización avanzada
% ------------------------------------------------------
\section{API y serialización avanzada}
\subsection{Validación avanzada y method validation}
\begin{lstlisting}[style=javastyle, caption={Groups y validación a nivel de método}]
import jakarta.validation.constraints.*;
import org.springframework.validation.annotation.Validated;

interface Create {}
interface Update {}

record Product(@NotNull(groups=Update.class) Long id,
               @NotBlank(groups={Create.class, Update.class}) String name) {}

@Validated
class ProductService {
  public @NotBlank String normalize(@NotBlank String name) { return name.trim(); }
}
\end{lstlisting}

\subsection{Jackson: configuración global y JavaTime}
\begin{lstlisting}[style=javastyle, caption={ObjectMapper personalizado}]
@Configuration
class JacksonCfg {
  @Bean com.fasterxml.jackson.databind.Module javaTime() {
    return new com.fasterxml.jackson.datatype.jsr310.JavaTimeModule();
  }
}
\end{lstlisting}

\subsection{HTTP caching, ETag y Problem+JSON}
\begin{lstlisting}[style=javastyle, caption={ETag condicional y Cache-Control}]
@GetMapping("/items/{id}")
ResponseEntity<ItemDto> get(@PathVariable Long id, WebRequest req) {
  var item = svc.get(id);
  var etag = '"' + item.version() + '"';
  if (req.checkNotModified(etag)) return null;
  return ResponseEntity.ok().eTag(etag).cacheControl(CacheControl.maxAge(Duration.ofMinutes(5))).body(item);
}
\end{lstlisting}

\begin{tipbox}
Estandariza errores con application/problem+json; usa un \texttt{@ControllerAdvice} que serialice código, título y detalle de la incidencia.
\end{tipbox}

% ------------------------------------------------------
% JPA avanzada
% ------------------------------------------------------
\section{JPA avanzada: concurrencia, caché y paginación}
\subsection{Optimistic/Pessimistic locking}
\begin{lstlisting}[style=javastyle, caption={@Version y bloqueo pesimista}] 
@Entity
class Account {
  @Id Long id;
  @Version Long version;
  BigDecimal balance;
}

interface AccountRepo extends JpaRepository<Account, Long> {
  @Lock(LockModeType.PESSIMISTIC_WRITE)
  @Query("select a from Account a where a.id=:id")
  Account findForUpdate(@Param("id") Long id);
}
\end{lstlisting}

\subsection{Caché de segundo nivel y batch}
\begin{lstlisting}[style=yamlstyle, caption={Ehcache/Hibernate y batch}] 
spring:
  jpa:
    properties:
      hibernate.cache.use_second_level_cache: true
      hibernate.cache.region.factory_class: org.hibernate.cache.jcache.JCacheRegionFactory
      hibernate.javax.cache.provider: org.ehcache.jsr107.EhcacheCachingProvider
      hibernate.default_batch_fetch_size: 50
\end{lstlisting}

\subsection{Paginación correcta}
\begin{lstlisting}[style=javastyle, caption={Page vs Slice y orden estable}] 
Page<User> page = repo.findAll(PageRequest.of(0, 50, Sort.by("id")));
Slice<User> slice = repo.findAllByActiveTrue(PageRequest.of(0, 50));
\end{lstlisting}

% ------------------------------------------------------
% Mapping y DTOs
% ------------------------------------------------------
\section{Mapeo eficiente con MapStruct}
\begin{lstlisting}[style=javastyle, caption={Mapper con @Mapper(componentModel="spring")}] 
@Mapper(componentModel = "spring")
interface UserMapper {
  UserDto toDto(User e);
  @Mapping(target="id", ignore=true)
  User toEntity(CreateUserDto dto);
}
\end{lstlisting}

% ------------------------------------------------------
% Seguridad OAuth2/OIDC avanzada
% ------------------------------------------------------
\section{OAuth2/OIDC avanzado}
\begin{lstlisting}[style=javastyle, caption={Client OIDC (login) y Resource Server}]
@Configuration
class OAuthCfg {
  @Bean SecurityFilterChain login(HttpSecurity http) throws Exception {
    http.securityMatcher("/ui/**")
      .authorizeHttpRequests(a -> a.anyRequest().authenticated())
      .oauth2Login(Customizer.withDefaults());
    return http.build();
  }
  @Bean SecurityFilterChain api(HttpSecurity http) throws Exception {
    http.securityMatcher("/api/**")
      .authorizeHttpRequests(a -> a.anyRequest().authenticated())
      .oauth2ResourceServer(o -> o.jwt());
    return http.build();
  }
}
\end{lstlisting}

% ------------------------------------------------------
% Cloud/Kubernetes
% ------------------------------------------------------
\section{Despliegue cloud-native y Kubernetes}
\begin{lstlisting}[style=yamlstyle, caption={Probes y recursos}] 
apiVersion: apps/v1
kind: Deployment
spec:
  template:
    spec:
      containers:
      - name: app
        image: app:latest
        readinessProbe:
          httpGet: { path: /actuator/health/readiness, port: 8080 }
        livenessProbe:
          httpGet: { path: /actuator/health/liveness, port: 8080 }
        resources:
          requests: { cpu: "200m", memory: "256Mi" }
          limits:   { cpu: "1",    memory: "512Mi" }
\end{lstlisting}

% ------------------------------------------------------
% OpenAPI/Swagger
% ------------------------------------------------------
\section{OpenAPI y documentación}
\begin{lstlisting}[style=javastyle, caption={springdoc-openapi configuración mínima}] 
// build: añade dependency org.springdoc:springdoc-openapi-starter-webmvc-ui
@Configuration
class OpenApiCfg {}
// UI: /swagger-ui.html  Spec: /v3/api-docs
\end{lstlisting}

% ------------------------------------------------------
% Logging estructurado
% ------------------------------------------------------
\section{Logging estructurado y correlación}
\begin{lstlisting}[style=javastyle, caption={MDC para correlation id}]
import org.slf4j.MDC;

try (var ignored = MDC.putCloseable("corrId", UUID.randomUUID().toString())) {
  log.info("procesando");
}
\end{lstlisting}

\begin{lstlisting}[style=yamlstyle, caption={Logback JSON (ejemplo)}]
logging:
  pattern:
    console: "%d{ISO8601} %level %logger - %X{corrId} %msg%n"
\end{lstlisting}

% ------------------------------------------------------
% Build avanzado
% ------------------------------------------------------
\section{Build: Gradle y Jib}
\begin{lstlisting}[style=bashstyle, caption={Gradle y Jib}]
./gradlew bootRun
./gradlew bootJar
./gradlew jibDockerBuild # com.google.cloud.tools:jib-spring-boot-extension
\end{lstlisting}

% ------------------------------------------------------
% Reactivo con R2DBC
% ------------------------------------------------------
\section{R2DBC: datos reactivos}
\begin{lstlisting}[style=javastyle, caption={Repositorio reactive y tx}] 
interface UserRxRepo extends org.springframework.data.r2dbc.repository.R2dbcRepository<User, Long> {}

@Service
class RxSvc {
  private final org.springframework.transaction.reactive.TransactionalOperator tx;
  private final UserRxRepo repo;
  Mono<User> create(User u) {
    return tx.execute(status -> repo.save(u));
  }
}
\end{lstlisting}

% ------------------------------------------------------
% Patrón Outbox
% ------------------------------------------------------
\section{Patrón Outbox para consistencia}
\begin{lstlisting}[style=javastyle, caption={Entidad Outbox y publicador}]
@Entity class OutboxEvent { @Id @GeneratedValue Long id; String type; String payload; Instant createdAt; }

@Service
class OutboxPublisher {
  @Transactional
  public void createAndPublish(Order order) {
    // 1) Persistir dominio
    // 2) Persistir evento outbox en la misma TX
  }
}
\end{lstlisting}

\appendix
\section{Comandos CLI y propiedades clave}
\begin{itemize}[leftmargin=*]
  \item \texttt{mvn spring-boot:run} \quad \texttt{./mvnw}\, disponible en proyectos generados.
  \item \texttt{mvn -DskipTests package} \quad empaqueta JAR.
  \item \texttt{mvn spring-boot:build-image} \quad imagen OCI con Buildpacks.
  \item Perfiles: \texttt{--spring.profiles.active=prod}. Variables: \texttt{SPRING\_APPLICATION\_NAME}.
\end{itemize}

\section{Cadena de filtros y flujo MVC (internals)}
\begin{itemize}[leftmargin=*]
  \item Servlet Container \textrightarrow{} FilterChain \textrightarrow{} DispatcherServlet \textrightarrow{} HandlerMapping \textrightarrow{} HandlerAdapter \textrightarrow{} Controller/Handler \textrightarrow{} ReturnValueHandlers \textrightarrow{} ViewResolver.
  \item Transacciones y proxies: \texttt{@Transactional} crea \emph{proxies} (JDK/CGLIB); fuga si se llama internamente al método desde la misma instancia.
  \item Post-procesadores: \texttt{BeanFactoryPostProcessor} y \texttt{BeanPostProcessor} permiten modificar definiciones/instancias en el ciclo de vida.
\end{itemize}

% ------------------------------------------------------
% Certificación Java SE 17 Developer (OCP)
% ------------------------------------------------------
\section{Certificación Java SE 17 Developer (OCP) \textemdash{} Guía práctica}
Esta guía resume los objetivos clave de OCP Java 17 (1Z0-829) y añade ejemplos y ejercicios.

\subsection{Temario esencial}
\begin{itemize}[leftmargin=*]
  \item Lenguaje: tipos, operadores, control de flujo, alcance, \texttt{var}, \texttt{record}, \texttt{sealed}, \texttt{switch expression}.
  \item POO: herencia, interfaces, clases anidadas, inmutabilidad, \texttt{equals/hashCode/compareTo}.
  \item Genéricos y colecciones: listas, conjuntos, mapas, comparadores, utilidades de \texttt{Collections}.
  \item Streams y funcional: pipelines, \texttt{Optional}, \texttt{Collectors}.
  \item Concurrencia: hilos, \texttt{ExecutorService}, colecciones concurrentes, \texttt{CompletableFuture}.
  \item Excepciones: jerarquía, \texttt{try-with-resources}, multi-\texttt{catch}.
  \item I/O y NIO.2: \texttt{Path}, \texttt{Files}, \texttt{WatchService}.
  \item Fecha/hora y i18n: \texttt{java.time}, \texttt{Locale}, \texttt{ResourceBundle}.
  \item Módulos (JPMS): \texttt{module-info.java}.
  \item JVM y herramientas: empaquetado, JAR, JFR/JCMD básico.
\end{itemize}

\subsection{Novedades relevantes y sintaxis moderna}
\begin{lstlisting}[style=javastyle, caption={Record, pattern matching y switch expression}]
public record Point(int x, int y) {}

static String typeOf(Object o) {
  if (o instanceof String s) return "str:" + s.length();
  if (o instanceof Number n) return "num:" + n.intValue();
  return "other";
}

static int dayToNum(String d) {
  return switch (d) {
    case "MON", "TUE", "WED", "THU", "FRI" -> 1;
    case "SAT", "SUN" -> 2;
    default -> throw new IllegalArgumentException();
  };
}
\end{lstlisting}

\begin{lstlisting}[style=javastyle, caption={Clases selladas (sealed) y jerarquías cerradas}]
public sealed interface Shape permits Circle, Rect {}
public final class Circle implements Shape {}
public non-sealed class Rect implements Shape {}
\end{lstlisting}

\subsection{POO sólida}
\begin{lstlisting}[style=javastyle, caption={equals/hashCode/compareTo correctos}]
public final class User implements Comparable<User> {
  private final String id; private final String email;
  public User(String id, String email) { this.id = id; this.email = email; }
  @Override public boolean equals(Object o) {
    if (this == o) return true; if (!(o instanceof User u)) return false;
    return id.equals(u.id);
  }
  @Override public int hashCode() { return id.hashCode(); }
  @Override public int compareTo(User o) { return this.id.compareTo(o.id); }
}
\end{lstlisting}

\subsection{Genéricos y colecciones}
\begin{lstlisting}[style=javastyle, caption={PECS: Producer Extends, Consumer Super}]
static <T> void copy(List<? extends T> src, List<? super T> dst) {
  for (T t : src) dst.add(t);
}
\end{lstlisting}

\begin{lstlisting}[style=javastyle, caption={Operaciones de Map y ordenación}]
var m = new java.util.HashMap<String,Integer>();
m.merge("a",1,Integer::sum);
m.merge("a",1,Integer::sum); // {a=2}
var sorted = m.entrySet().stream()
  .sorted(java.util.Map.Entry.<String,Integer>comparingByValue().reversed())
  .toList();
\end{lstlisting}

\subsection{Streams, Optional y Collectors}
\begin{lstlisting}[style=javastyle, caption={Pipeline típico y colecciones inmutables}]
var list = java.util.List.of(1,2,3,4,5);
var evensSquared = list.stream()
  .filter(i -> i % 2 == 0)
  .map(i -> i * i)
  .collect(java.util.stream.Collectors.toUnmodifiableList());
\end{lstlisting}

\begin{lstlisting}[style=javastyle, caption={Collectors avanzados}]
var byParity = list.stream().collect(java.util.stream.Collectors.partitioningBy(i -> i % 2 == 0));
var stats = list.stream().collect(java.util.stream.Collectors.summarizingInt(Integer::intValue));
\end{lstlisting}

\subsection{Concurrencia y paralelismo}
\begin{lstlisting}[style=javastyle, caption={ExecutorService y CompletableFuture}]
var pool = java.util.concurrent.Executors.newFixedThreadPool(4);
try {
  var f = java.util.concurrent.CompletableFuture.supplyAsync(() -> 40 + 2, pool)
    .thenApply(x -> x * 2)
    .orTimeout(1, java.util.concurrent.TimeUnit.SECONDS);
  System.out.println(f.join());
} finally { pool.shutdown(); }
\end{lstlisting}

\subsection{Excepciones y TWR}
\begin{lstlisting}[style=javastyle, caption={Try-with-resources y supresión}]
try (var in = java.nio.file.Files.newBufferedReader(java.nio.file.Path.of("data.txt"))) {
  System.out.println(in.readLine());
} catch (java.io.IOException e) {
  e.printStackTrace();
}
\end{lstlisting}

\subsection{I/O, NIO.2 e i18n}
\begin{lstlisting}[style=javastyle, caption={Path/Files y WatchService}]
var path = java.nio.file.Path.of("/tmp");
try (var s = java.nio.file.Files.walk(path)) {
  s.filter(java.nio.file.Files::isRegularFile).forEach(System.out::println);
}
var ws = java.nio.file.FileSystems.getDefault().newWatchService();
path.register(ws, java.nio.file.StandardWatchEventKinds.ENTRY_CREATE);
\end{lstlisting}

\begin{lstlisting}[style=javastyle, caption={java.time y ResourceBundle}]
var now = java.time.ZonedDateTime.now();
var fmt = java.time.format.DateTimeFormatter.ofPattern("yyyy-MM-dd HH:mm z", java.util.Locale.US);
System.out.println(now.format(fmt));
var b = java.util.ResourceBundle.getBundle("messages", new java.util.Locale("es","ES"));
System.out.println(b.getString("hello"));
\end{lstlisting}

\subsection{JPMS y herramientas}
\begin{lstlisting}[style=javastyle, caption={module-info.java mínimo}]
module com.example.app {
  requires java.logging;
  exports com.example.api;
}
\end{lstlisting}

\subsection{Trampas típicas y ejercicios}
\begin{lstlisting}[style=javastyle, caption={Autoboxing y cache de Integer}]
Integer a = 128, b = 128; System.out.println(a == b); // false
Integer x = 127, y = 127; System.out.println(x == y); // true
\end{lstlisting}

\begin{itemize}[leftmargin=*]
  \item Identifica salida de código: orden de inicialización, sobrecarga vs sobreescritura.
  \item Corrige genéricos: comodines y límites múltiples; PECS.
  \item Completa TWR, \texttt{Files.walk} y \texttt{Collectors.groupingBy}.
\end{itemize}

\subsection{Plan de estudio (8 semanas)}
\begin{enumerate}[leftmargin=*]
  \item Sintaxis base y POO.  \item Genéricos y colecciones.
  \item Streams/Optional. \item Excepciones, I/O y NIO.2.
  \item Concurrencia: hilos, locks, ExecutorService, CompletableFuture.
  \item Fecha/hora e i18n; módulos JPMS.
  \item Simulacros de examen y repasos.
  \item JFR/JCMD básico y último simulacro.
\end{enumerate}

\begin{tipbox}
Practica con bancos de preguntas y escribe micro-programas que demuestren cada objetivo. Prioriza la lectura de preguntas con trampas de tipos y promoción numérica.
\end{tipbox}

\end{document}